%!TEX program = xelatex
\PassOptionsToPackage{full}{textcomp}
\documentclass[lang=en,11pt,a4paper,citestyle=authoryear]{elegantpaper}

\title{HW5 - STAT 709}
\author{Wells Guan}

\usepackage{array,url,stix}
\usepackage{tikz}
\usepackage{tikz-cd}
\usepackage{enumitem}
\usepackage{amsmath}
\usepackage{amssymb}
\usepackage{listings}
\usepackage{subfigure}

\lstset{
    frame=single,
    columns=fullflexible,
    keepspaces=true,
    showstringspaces=false,
    breaklines=true,
    breakatwhitespace=false,
    basicstyle=\ttfamily\footnotesize,
    numbers=left,
    numberstyle=\tiny,
    numbersep=6pt
}

\newcommand{\Z}{\mathbb{Z}}
\newcommand{\R}{\mathbb{R}}
\newcommand{\N}{\mathbb{N}}
\newcommand{\B}{\mathcal{B}}
\newcommand{\E}{\mathbb{E}}
\renewcommand{\P}{\mathbb{P}}
\newcommand{\norm}[1]{\left\lVert#1\right\rVert}

\begin{document}
\maketitle

\noindent All exercises are required unless marked optional.

\subsection*{Problem 1. Empirical distributions and moment estimates}
Bickel and Doksum, \textit{Mathematical Statistics: Basic Ideas and
Selected Topics}, Volume I, second edition, Section 2.5,
Problems for Section 2.1: Exercises 6 and 9 (required).

\subsubsection*{Exercise 2.1.6}
Let $X_{(1)}\le\cdots\le X_{(n)}$ be the order statistics of a
sample $X_1,\ldots,X_n$ (see Problem B.2.8). There is a one-to-one
correspondence between the empirical distribution function $\widehat F$
and the order statistics in the sense that, given the order statistics,
we may construct $\widehat F$, and given $\widehat F$, we know the
order statistics. Give the details of this correspondence.

\par\vspace{0.5em}\noindent\textbf{Sol.}
For the fixed sample size $n$, the empirical distribution function is
\[
\widehat F(x)=\frac1n\sum_{i=1}^n\mathbf1_{\{X_i\le x\}}.
\]
Reordering a sample does not change this sum. Thus the order statistics
determine the entire function through
\[
{\widehat F(x)
=\frac1n\sum_{i=1}^n\mathbf1_{\{X_{(i)}\le x\}},\qquad x\in\R.}
\]
Conversely, for every $j=1,\ldots,n$,
\[
{X_{(j)}=\inf\{x\in\R:\widehat F(x)\ge j/n\}}
\]
and we are done.

\subsubsection*{Exercise 2.1.9}
Suppose $\mathbf X=(X_1,\ldots,X_n)$, where the $X_i$ are independent
$N(0,\sigma^2)$.
\begin{enumerate}[label=(\alph*),leftmargin=*]
\item Find an estimate of $\sigma^2$ based on the second moment.
\item Construct an estimate of $\sigma$ using the estimate of part (a)
and the equation $\sigma=\sqrt{\sigma^2}$.
\item Use the empirical substitution principle to construct an estimate
of $\sigma$ using the relation $\E|X_1|=\sigma\sqrt{2/\pi}$.
\end{enumerate}

\par\vspace{0.5em}\noindent\textbf{Sol.}
The empirical substitution principle replaces an expectation under the
unknown distribution by its empirical average.
\begin{enumerate}[label=(\alph*),leftmargin=*]
\item Since $\E X_1^2=\sigma^2$, replacing the second moment by its
sample analogue gives
\[
{\widehat{\sigma^2}=\frac1n\sum_{i=1}^nX_i^2}
\]
with $\E\widehat{\sigma^2}=\sigma^2$.

\item Applying the square-root relation gives
\[
{\widehat\sigma_{\mathrm{rms}}
=\sqrt{\frac1n\sum_{i=1}^nX_i^2}.}
\]


\item Rewriting the given relation as
$\sigma=\sqrt{\pi/2}\,\E|X_1|$ gives the alternative estimate
\[
{\widehat\sigma_{\mathrm{abs}}
=\sqrt{\frac\pi2}\,\frac1n\sum_{i=1}^n|X_i|.}
\]
In particular, this estimate is unbiased. For completeness, the relation
follows directly from the normal density:
\[
\E|X_1|
=\frac{2}{\sigma\sqrt{2\pi}}
\int_0^\infty x e^{-x^2/(2\sigma^2)}\,dx
=\sigma\sqrt{\frac2\pi}.
\]
Both estimates of $\sigma$ are consistent by the law of large numbers
and continuity of the square-root map.
\end{enumerate}

\subsection*{Problem 2. Maximum likelihood estimates}
Bickel and Doksum, the same edition, Section 2.5,
Problems for Section 2.2: Exercises 12 and 13 (required).

\subsubsection*{Exercise 2.2.12}
Let $X_1,\ldots,X_n$ be i.i.d.\ as $X\sim P_\theta$, with
$\theta\in\Theta\subseteq\R^d$ identifiable. Suppose $X$ has possible
values $v_1,\ldots,v_k$ and that
\[
q(\theta)=h(\mu_1(\theta),\ldots,\mu_r(\theta))
\]
for some $\R^k$-valued function $h$. Show that the method-of-moments
estimate $\widehat q=h(\widehat\mu_1,\ldots,\widehat\mu_r)$ can be
written as a frequency plug-in estimate.

\par\vspace{0.5em}\noindent\textbf{Sol.}
Write $p_\ell(\theta)=\P_\theta(X=v_\ell)$ and let
\[
\widehat p_\ell=\frac1n\sum_{i=1}^n\mathbf1_{\{X_i=v_\ell\}},
\qquad \ell=1,\ldots,k,
\]
be the empirical relative frequencies. The theoretical and empirical
raw moments satisfy
\[
\mu_j(\theta)=\sum_{\ell=1}^k v_\ell^j p_\ell(\theta),
\qquad
\widehat\mu_j=\frac1n\sum_{i=1}^nX_i^j
=\sum_{\ell=1}^k v_\ell^j\widehat p_\ell.
\]
Define the functional of a probability vector $p$ by
\[
Q(p)=h\left(\sum_{\ell=1}^k v_\ell p_\ell,\ldots,
\sum_{\ell=1}^k v_\ell^r p_\ell\right).
\]
Then $q(\theta)=Q(p(\theta))$, and replacing the probabilities by their
empirical frequencies yields
\[
{Q(\widehat p)
=h(\widehat\mu_1,\ldots,\widehat\mu_r)=\widehat q.}
\]
Thus the moment estimate is exactly a frequency plug-in estimate,
not merely an asymptotic approximation. As usual, the displayed
estimator is understood on the sample values for which $h$ is defined.

\subsubsection*{Exercise 2.2.13}
\textit{General method of moment estimates.}
Suppose $X_1,\ldots,X_n$ are i.i.d.\ as $X\sim P_\theta$, with
$\theta\in\Theta\subseteq\R^d$ identifiable. Let $g_1,\ldots,g_r$ be
given linearly independent functions and write
\[
\mu_j(\theta)=\E_\theta g_j(X),\qquad
\widehat\mu_j=\frac1n\sum_{i=1}^n g_j(X_i),
\qquad j=1,\ldots,r.
\]
Suppose $X$ has possible values $v_1,\ldots,v_k$ and that
\[
q(\theta)=h(\mu_1(\theta),\ldots,\mu_r(\theta))
\]
for some $\R^k$-valued function $h$.
\begin{enumerate}[label=(\alph*),leftmargin=*]
\item Show that the method-of-moments estimate
$\widehat q=h(\widehat\mu_1,\ldots,\widehat\mu_r)$ is a frequency
plug-in estimate.
\item Suppose $\{P_\theta:\theta\in\Theta\}$ is the $k$-parameter
exponential family given by (1.6.10). Let $g_j(X)=T_j(X)$,
$1\le j\le k$. In the following cases, find the method-of-moments
estimates:
\begin{enumerate}[label=(\roman*),leftmargin=*]
\item Beta, $\beta(1,\theta)$;
\item Beta, $\beta(\theta,1)$;
\item Rayleigh,
$p(x,\theta)=(x/\theta^2)\exp(-x^2/(2\theta^2))$,
$x>0$, $\theta>0$;
\item Gamma, $\Gamma(p,\theta)$, $p$ fixed;
\item Inverse Gaussian, $IG(\mu,\lambda)$, $\theta=(\mu,\lambda)$.
See Problem 1.6.36.
\end{enumerate}
\textit{Hint:} Use Corollary 1.6.1.
\end{enumerate}

\par\vspace{0.5em}\noindent\textbf{Sol.}
\begin{enumerate}[label=(\alph*),leftmargin=*]
\item With $p_\ell(\theta)$ and $\widehat p_\ell$ as in Exercise 2.2.12,
\[
\mu_j(\theta)=\sum_{\ell=1}^k g_j(v_\ell)p_\ell(\theta),
\qquad
\widehat\mu_j=\sum_{\ell=1}^k g_j(v_\ell)\widehat p_\ell.
\]
Consequently, if
\[
Q(p)=h\left(\sum_{\ell=1}^k g_1(v_\ell)p_\ell,\ldots,
\sum_{\ell=1}^k g_r(v_\ell)p_\ell\right),
\]
then
\[
{\widehat q=Q(\widehat p).}
\]
This is the same substitution argument as before, with $x^j$
replaced by the prescribed functions $g_j(x)$.

\item Here the moment functions are the exponential-family canonical
statistics, rather than necessarily $X,X^2,\ldots$.
For a density written as
\[
f_\eta(x)=a(x)\exp\{\eta^\top T(x)-A(\eta)\},
\]
the moment identity is $\E_\eta T(X)=\nabla A(\eta)$.
Equating this expectation to the empirical mean of $T(X)$ gives the
required generalized moment equations. The finite-support assumption
in part (a) is not imposed on these continuous examples.

\begin{enumerate}[label=(\roman*),leftmargin=*]
\item \textbf{Beta $(1,\theta)$.}
The density is
\[
f_\theta(x)=\theta(1-x)^{\theta-1},\qquad0<x<1,\quad\theta>0.
\]
Its canonical statistic is $T(x)=\log(1-x)$ and
\[
\E_\theta\log(1-X)=-\frac1\theta.
\]
For example, this follows by setting $y=1-x$ and differentiating
$\int_0^1y^{\theta-1}\,dy=1/\theta$ with respect to $\theta$.
Matching the empirical logarithmic moment gives
\[
{\widehat\theta
=-\left(\frac1n\sum_{i=1}^n\log(1-X_i)\right)^{-1}
=-\frac{n}{\sum_{i=1}^n\log(1-X_i)}.}
\]
Since $0<X_i<1$ almost surely, the denominator is strictly negative.
The raw-first-moment estimate $\overline X^{-1}-1$ is a different
moment estimate and is not the one requested when $g=T$.

\item \textbf{Beta $(\theta,1)$.}
Now
\[
f_\theta(x)=\theta x^{\theta-1},\qquad0<x<1,
\]
so $T(x)=\log x$ and $\E_\theta\log X=-1/\theta$.
Thus
\[
{\widehat\theta
=-\left(\frac1n\sum_{i=1}^n\log X_i\right)^{-1}
=-\frac{n}{\sum_{i=1}^n\log X_i}.}
\]
Again this is based on the specified canonical statistic;
the raw-first-moment estimate would instead be
$\overline X/(1-\overline X)$.

\item \textbf{Rayleigh.}
The canonical statistic is $T(x)=x^2$ because the parameter-dependent
part of the log density is $-2\log\theta-x^2/(2\theta^2)$.
The variable $X^2/(2\theta^2)$ has the standard exponential
distribution, so $\E_\theta X^2=2\theta^2$.
Consequently,
\[
{\widehat\theta
=\sqrt{\frac1{2n}\sum_{i=1}^nX_i^2}.}
\]

\item \textbf{Gamma with known shape $p$.}
To specify the parameter convention explicitly, first use the
shape--scale density
\[
f_\theta(x)=\frac{x^{p-1}}{\Gamma(p)\theta^p}e^{-x/\theta},
\qquad x>0,\quad\theta>0.
\]
Then $T(x)=x$ and $\E_\theta X=p\theta$, whence
\[
{\widehat\theta=\overline X/p\quad\text{(scale convention)}.}
\]
If $\Gamma(p,\theta)$ is instead defined with $\theta$ as the
\emph{rate}, namely
$f_\theta(x)=\theta^p x^{p-1}e^{-\theta x}/\Gamma(p)$,
then $\E_\theta X=p/\theta$ and the corresponding answer is
\[
{\widehat\theta=p/\overline X\quad\text{(rate convention)}.}
\]
The question screenshot does not include the textbook's earlier
Gamma density definition, so both conventions are stated explicitly.

\item \textbf{Inverse Gaussian.}
Use the mean--shape parameterization
\[
f_{\mu,\lambda}(x)
=\left(\frac{\lambda}{2\pi x^3}\right)^{1/2}
\exp\left[-\frac{\lambda(x-\mu)^2}{2\mu^2x}\right],
\qquad x>0,\quad\mu,\lambda>0.
\]
Expanding the exponent gives
\[
-\frac{\lambda(x-\mu)^2}{2\mu^2x}
=-\frac{\lambda}{2\mu^2}x-\frac\lambda2\frac1x
+\frac\lambda\mu.
\]
Hence take $T_1(x)=x$, $T_2(x)=1/x$.
Their expectations are
\[
\E X=\mu,\qquad \E\frac1X=\frac1\mu+\frac1\lambda.
\]
One way to verify both identities is to use the normalizing integral
\[
Z(a,b)=\int_0^\infty x^{-3/2}e^{-ax-b/x}\,dx
=\sqrt{\frac\pi b}\,e^{-2\sqrt{ab}},\qquad a,b>0.
\]
The identity follows from normalization of the displayed inverse
Gaussian density with $a=\lambda/(2\mu^2)$ and $b=\lambda/2$;
every $a,b>0$ arises from such a parameter pair.
Differentiating the logarithm of this integral yields
\[
\E X=-\partial_a\log Z=\sqrt{b/a}=\mu,
\qquad
\E\frac1X=-\partial_b\log Z
=\frac1{2b}+\sqrt{a/b}=\frac1\lambda+\frac1\mu.
\]
The differentiation is justified locally by the exponential decay at
both zero and infinity.
Let $H_n=n^{-1}\sum_iX_i^{-1}$. The moment equations are
\[
\overline X=\widehat\mu,
\qquad H_n=\frac1{\widehat\mu}+\frac1{\widehat\lambda}.
\]
Their solution is
\[
{\widehat\mu=\overline X,\qquad
\widehat\lambda=\left(H_n-\frac1{\overline X}\right)^{-1}.}
\]
By Cauchy--Schwarz,
\[
\left(\frac1n\sum_iX_i\right)
\left(\frac1n\sum_iX_i^{-1}\right)\ge1,
\]
with equality exactly when all observations coincide. Thus the
denominator is positive almost surely for $n\ge2$.
For $n=1$ (or a constant sample), there is no finite solution for
$\lambda$; the formal estimate lies at $\lambda=\infty$.

For comparison only, matching the ordinary moments $X$ and $X^2$
would instead give $\widehat\mu=\overline X$ and
$\widehat\lambda=\overline X^3/
[n^{-1}\sum_i(X_i-\overline X)^2]$, using
$\operatorname{Var}(X)=\mu^3/\lambda$.
This is not the estimator requested by the choice $g_j=T_j$.
\end{enumerate}
\end{enumerate}

\subsection*{Problem 3. Method of moments versus maximum likelihood}
Let
\[
X_1,\ldots,X_n\overset{\mathrm{i.i.d.}}{\sim}
\operatorname{Unif}([\theta_1,\theta_2]),
\qquad n\ge2,\quad\theta_1<\theta_2,
\]
where both endpoints are unknown. Write
\[
\overline X=\frac1n\sum_{i=1}^nX_i,
\qquad
S^2=\frac1{n-1}\sum_{i=1}^n(X_i-\overline X)^2.
\]
A method-of-moments estimator is
\[
\widehat\theta_1
=\overline X-\sqrt{\frac{3(n-1)}n S^2},
\qquad
\widehat\theta_2
=\overline X+\sqrt{\frac{3(n-1)}n S^2}.
\]
\begin{enumerate}[label=(\roman*),leftmargin=*]
\item Find the maximum likelihood estimators of $\theta_1$ and $\theta_2$.
\item Compare the mean squared errors (MSEs) of the
method-of-moments estimator and the maximum likelihood estimator.
Which estimator is better?
\end{enumerate}

\par\vspace{0.5em}
\noindent\textbf{Sol.}
To distinguish the two procedures, retain
$(\widehat\theta_1,\widehat\theta_2)$ for the method-of-moments
estimator and write
$(\widetilde\theta_1,\widetilde\theta_2)$ for the MLE.

\begin{enumerate}[label=(\roman*),leftmargin=*]
\item Let $X_{(1)}=\min_iX_i$ and $X_{(n)}=\max_iX_i$.
The likelihood is
\[
L(\theta_1,\theta_2;x)
=(\theta_2-\theta_1)^{-n}
\mathbf1_{\{\theta_1\le x_{(1)},\ x_{(n)}\le\theta_2\}},
\qquad\theta_1<\theta_2.
\]
For positive likelihood, the proposed interval must contain all the
observations. Over these intervals, the likelihood is strictly decreasing
in their length. If $x_{(1)}<x_{(n)}$, the shortest admissible interval is
$[x_{(1)},x_{(n)}]$. Therefore
\[
{\widetilde\theta_1=X_{(1)},\qquad
\widetilde\theta_2=X_{(n)}.}
\]
Since the distribution is continuous and $n\ge2$, the observations
are not all equal almost surely, so this MLE exists and is unique
almost surely. On the null event that all observations coincide,
the likelihood is unbounded as the interval length tends to zero.

\item Put $d=\theta_2-\theta_1>0$ and standardize by
\[
U_i=\frac{X_i-\theta_1}{d}
\overset{\mathrm{i.i.d.}}{\sim}\operatorname{Unif}([0,1]).
\]
Both procedures are location--scale equivariant, so their endpoint
MSEs are $d^2$ times the corresponding MSEs on $[0,1]$.
We first compare each endpoint separately. Under squared Euclidean
loss for the endpoint vector, the risk is the sum of these two MSEs.

\medskip
\noindent\textbf{Exact MSEs of the MLE.}
For $0\le u\le1$,
\[
\P(U_{(1)}>u)=(1-u)^n,
\]
so $U_{(1)}$ has density $n(1-u)^{n-1}$ and
\[
\E U_{(1)}^2
=n\int_0^1u^2(1-u)^{n-1}\,du
=\frac{2}{(n+1)(n+2)}.
\]
By reflection, $1-U_{(n)}$ has the same distribution as $U_{(1)}$.
Consequently,
\[
{
\operatorname{MSE}(\widetilde\theta_1)
=\operatorname{MSE}(\widetilde\theta_2)
=\frac{2d^2}{(n+1)(n+2)}.}
\]
In particular, the total squared-error risk of the MLE is
$4d^2/[(n+1)(n+2)]$.

\medskip
\noindent\textbf{An exact expression for the moment-estimator MSEs.}
Write
\[
\overline U=\frac1n\sum_iU_i,
\qquad V_n=\frac1n\sum_i(U_i-\overline U)^2,
\qquad A_n=\E\sqrt{V_n}.
\]
The lower and upper moment estimates on $[0,1]$ are
$\overline U-\sqrt{3V_n}$ and $\overline U+\sqrt{3V_n}$.
The transformation $U_i\mapsto1-U_i$ preserves the joint law,
sends $\overline U$ to $1-\overline U$, and leaves $V_n$ unchanged.
It follows that
\[
\E[\overline U\sqrt{V_n}]
=\E[(1-\overline U)\sqrt{V_n}]
=\frac12 A_n.
\]
Also,
\[
\E\overline U^2=\frac14+\frac1{12n},
\qquad \E V_n=\frac{n-1}{12n}.
\]
Thus the two endpoint MSEs are equal, and their exact expression is
\[
{
\operatorname{MSE}(\widehat\theta_j)
=d^2\left(\frac12-\frac1{6n}-\sqrt3\,A_n\right),
\qquad j=1,2.}
\]
Here the expectation is entirely parameter-free; explicitly,
\[
A_n=\int_{[0,1]^n}
\sqrt{\frac1n\sum_{i=1}^n(u_i-\overline u)^2}
\,du_1\cdots du_n,
\qquad \overline u=\frac1n\sum_i u_i.
\]
There is no need to replace this expectation by the square root
of $\E V_n$, which would give an incorrect finite-sample MSE.
For any fixed $n$, the moment estimator has smaller endpoint MSE
exactly when
\[
A_n>
\frac1{\sqrt3}\left(
\frac12-\frac1{6n}-\frac{2}{(n+1)(n+2)}
\right).
\]

\medskip
\noindent\textbf{Small-sample comparison: $n=2$.}
In this case $\sqrt{V_2}=|U_1-U_2|/2$, and
\[
\E|U_1-U_2|
=2\int_0^1\int_0^v(v-u)\,du\,dv=\frac13.
\]
Therefore $A_2=1/6$, and
\[
\operatorname{MSE}(\widehat\theta_j)
=d^2\frac{5-2\sqrt3}{12}
<\frac{d^2}{6}
=\operatorname{MSE}(\widetilde\theta_j).
\]
Thus the moment estimator is better at $n=2$.

\medskip
\noindent\textbf{Large-sample comparison.}
Let $W_i=U_i-1/2$. For a uniform observation,
\[
\E W_i=0,\quad \E W_i^2=\frac1{12},\quad
\E W_i^3=0,\quad \E W_i^4=\frac1{80}.
\]
Hence
\[
\operatorname{Var}(W_i^2)=\frac1{80}-\frac1{144}
=\frac1{180},
\qquad \operatorname{Cov}(W_i,W_i^2)=0.
\]
Define $a_n=\overline W$ and
$b_n=n^{-1}\sum_i(W_i^2-1/12)$. Then
\[
V_n=\frac1{12}+b_n-a_n^2.
\]
The function $v\mapsto\sqrt{3v}$ has derivative $3$ at $v=1/12$.
Its expansion gives
\[
\begin{aligned}
\overline U-\sqrt{3V_n}&=a_n-3b_n+R_n,\\
\overline U+\sqrt{3V_n}-1&=a_n+3b_n-R_n,
\end{aligned}
\]
where $\E R_n^2=O(n^{-2})$.
To justify this remainder bound, on $|V_n-1/12|\le1/24$ the
square-root Taylor remainder is bounded by a constant times
$(V_n-1/12)^2$, and $R_n$ also includes $3a_n^2$.
Since the observations are bounded and centered,
\[
\E a_n^4=O(n^{-2}),\qquad
\E b_n^4=O(n^{-2}),\qquad
\E|V_n-1/12|^4=O(n^{-2}).
\]
The last identity follows from $V_n-1/12=b_n-a_n^2$
and $\E a_n^8=O(n^{-4})$.
The complementary event has probability $O(n^{-2})$ by Markov's
inequality; the remainder there is uniformly bounded.
Together these facts establish $\E R_n^2=O(n^{-2})$.

Finally,
\[
\E(a_n\pm3b_n)^2
=\frac1n\left(\frac1{12}+\frac9{180}\right)
=\frac{2}{15n}.
\]
Cauchy--Schwarz controls the cross term involving $R_n$, so
\[
{
\operatorname{MSE}(\widehat\theta_j)
=\frac{2d^2}{15n}+o(n^{-1}),
\qquad j=1,2.}
\]
In contrast,
\[
\operatorname{MSE}(\widetilde\theta_j)
=\frac{2d^2}{n^2}+o(n^{-2}).
\]
Therefore the MLE has smaller MSE for all sufficiently large $n$.
\end{enumerate}

\end{document}